

A consistent and well-defined computational environment was essential for ensuring reproducibility, stability, and fairness during the comparative evaluation of the four deep learning architectures. All experiments—including preprocessing, training, validation, and MLflow-based tracking—were executed on a single Apple Silicon system. Tables \ref{tab:hardware} and \ref{tab:software} summarise the key hardware and software configurations.

\section{Hardware Specifications}

\begin{table}[H]
\centering
\caption{Hardware Configuration Used for Model Training and Validation}
\label{tab:hardware}
\begin{tabular}{|l|l|}
\hline
\textbf{Component} & \textbf{Specification} \\ \hline
System & Apple MacBook (Apple Silicon) \\ \hline
Processor & Apple M4 \\ \hline
Unified Memory & 24 GB \\ \hline
GPU Backend & Apple GPU with Metal (MPS) \\ \hline
Storage & NVMe SSD \\ \hline
Operating System & macOS (latest version) \\ \hline
\end{tabular}
\end{table}

\section{Software Environment}

The complete machine learning workflow was executed using Python~3.11 and a curated set of libraries to ensure cross-framework compatibility and stable acceleration on Apple Silicon. The environment was managed using a pinned \texttt{requirements.txt} file.

\begin{table}[H]
\centering
\caption{Key Software and Library Versions}
\label{tab:software}
\begin{tabular}{|l|l|}
\hline
\textbf{Software Component} & \textbf{Version / Constraint} \\ \hline
Python & 3.11 \\ \hline
NumPy & $\geq$1.24.3, $<$1.26 \\ \hline
Pandas & 2.0.3 \\ \hline
SciPy & 1.11.3 \\ \hline
h5py & 3.9.0 \\ \hline
Scikit-Learn & 1.3.0 \\ \hline
TensorFlow & 2.15.0 \\ \hline
tensorflow-metal & 1.1.0 \\ \hline
PyTorch & 2.3.0 \\ \hline
Torchvision & 0.18.0 \\ \hline
Torchaudio & 2.3.0 \\ \hline
Matplotlib & 3.7.2 \\ \hline
Seaborn & 0.12.2 \\ \hline
Joblib & 1.3.2 \\ \hline
MLflow & 2.22.1 \\ \hline
psutil & $\geq$5.9.5, $<$6.0.0 \\ \hline
\end{tabular}
\end{table}

\section{Environment Summary}

The Apple M4 hardware, combined with Python~3.11 and a tightly controlled software stack, ensured consistent and reproducible execution of the training pipeline. Integration with MLflow further strengthened experiment governance, enabling robust comparison of all four candidate models under identical computational conditions.

\section{Project Source Code and Repository Structure}

All implementation work carried out for this dissertation is publicly available in a structured GitHub repository. 
The source code for the complete Phase~2 pipeline---including dataset loaders, deep neural network architectures, 
training scripts, MLOps automation utilities, and configuration files---is located at:

\begin{center}
\texttt{https://github.com/mohena-com/wifi\_project-phase-2/tree/dev\_1.0/src/phase2}
\end{center}

This directory forms the core of the project and contains all modules used for data handling, model design, 
experiment configuration, MLflow tracking, and inference. Table~\ref{tab:phase2files} summarises the 
key files within the \texttt{phase2} module along with their primary roles.

\begin{table}[H]
\centering
\caption{Summary of Important Source Files in the \texttt{phase2} Directory}
\label{tab:phase2files}
\begin{tabular}{|l|l|}
\hline
\textbf{File Name} & \textbf{Description} \\ \hline

CSI\_ID\_MLOPS\_advanced\_v1.3.py 
& Main MLOps-enabled training script with end-to-end automation \\ \hline

DL\_CSILSTMNet.py 
& Implementation of the CNN--LSTM hybrid model \\ \hline

DL\_DenseNet1D.py 
& Definition of the DenseNet1D architecture \\ \hline

DL\_EfficientNet1DLSTM.py 
& EfficientNet1D + LSTM combined architecture \\ \hline

DL\_MobileNetV3.py 
& Implementation of MobileNetV3 (1D adapted) \\ \hline

DS\_WifiCSIDataset.py 
& Custom dataset loader for Wi-Fi CSI frames and metadata \\ \hline

config\_reader.py 
& Utility for reading configuration properties and hyperparameters \\ \hline

csi\_id\_config.properties 
& Centralised configuration file for model, training, and dataset parameters \\ \hline



model\_definitions.py 
& Consolidated definitions of all neural network classes for ease of experimentation \\ \hline

\end{tabular}
\end{table}

The structure is intentionally modular to support iterative experimentation and systematic comparison of four 
deep learning architectures. Each model file is decoupled from dataset utilities and training logic, enabling 
plug-and-play experimentation. The primary MLOps script orchestrates configuration loading, model selection, 
MLflow logging, training execution, and evaluation, thereby enforcing reproducibility and consistency across 
all experimental runs.
